\subsection{When Backshift Is Not Necessary}

Backshift is common when the reporting verb is in the past, but it is not
always mandatory.

If the reported statement describes a fact that remains true, the original
tense may be retained.

\begin{quotation}
Direct: Sarah said, ``The Earth is round.''

Reported: Sarah said that the Earth is round.
\end{quotation}

The present tense is appropriate because the statement remains true.

Similarly:

\begin{quotation}
He says, ``I am tired.''
\end{quotation}

becomes:

\begin{quotation}
He says that he is tired.
\end{quotation}

When the reporting verb is in the present tense, backshift is generally
unnecessary.


\subsection{A Systematic Transformation Procedure}

A direct sentence can be converted into reported speech by following these
steps:

\begin{enumerate}
    \item Identify the reporting verb.
    \item Determine whether the original sentence is a statement, question,
    command, or request.
    \item Identify the referents of all pronouns.
    \item Apply the appropriate tense changes when backshift is required.
    \item Change relevant expressions of time and place.
    \item For questions, use statement word order.
    \item Check that the final sentence preserves the meaning of the original.
\end{enumerate}


\subsection{Complete Example}

Consider:

\begin{quotation}
David said to Maria, ``I will finish this project tomorrow because I am very
busy today.''
\end{quotation}

The transformation can be performed as follows:

\begin{enumerate}
    \item ``said to Maria'' can become ``told Maria''.
    \item ``I'' refers to David, so it becomes ``he''.
    \item ``will finish'' becomes ``would finish''.
    \item ``this'' becomes ``that''.
    \item ``tomorrow'' becomes ``the next day''.
    \item ``am'' becomes ``was''.
    \item ``today'' becomes ``that day''.
\end{enumerate}

The resulting sentence is:

\begin{quotation}
David told Maria that he would finish that project the next day because he
was very busy that day.
\end{quotation}


\subsection{Summary}

The principal transformations in reported speech are:

\begin{center}
\begin{tabular}{ll}
\textbf{Direct} & \textbf{Reported} \\ \hline
am/is/are & was/were \\
V1 & V2 \\
am/is/are + V-ing & was/were + V-ing \\
have/has + V3 & had + V3 \\
V2 & had + V3 \\
will & would \\
can & could \\
may & might
\end{tabular}
\end{center}

The fundamental structures are:

\[
\text{Statement: } \quad
\text{said/told + (that) + statement}
\]

\[
\text{Yes/No question: } \quad
\text{asked + if/whether + statement}
\]

\[
\text{WH-question: } \quad
\text{asked + WH-word + statement}
\]

\[
\text{Command/request: } \quad
\text{told/asked + object + to + V1}
\]

The transformation should not be treated as a mechanical substitution of
words. Pronouns, tense, time, and place must be changed according to the
relationship between the original statement and the context in which it is
reported.
